\section{Discussion}

\paragraph{Standard RAG reporting is insufficient.}
The central result is not that one retrieval policy should replace another.
It is that a common evidential pattern in RAG papers is too weak: local
retrieval score improves, one faithfulness metric changes, and a deployment
claim follows. Our controls show that each step can fail. Local relevance does
not identify answer support; CCS alone does not identify answer support; a
single automatic metric does not identify effect size; and one baseline row
does not identify a deployment winner.

\paragraph{Why the null result matters.}
The matched-coherence null is not a failed paper; it is the paper's strongest
control. Without it, one might conclude from correlations and pilot tables that
coherence is the mechanism behind refinement failures. With it, the evidence
supports a more careful interpretation: coherence is one retrieval-set property
that can diagnose fragility, but support also depends on answer coverage,
redundancy, contradiction, generator behavior, and evaluator bias. The collapse
of the intuitive story is the point: it shows why identifiability controls are
needed before RAG papers turn retrieval statistics into faithfulness claims.

\paragraph{What CCS is after the controls.}
CCS is a diagnostic variable, not a mechanism. It is useful because it measures
a set-level property that local relevance misses and because the noise
experiment shows that coherence-preserving degradation behaves differently
from random degradation. It should not be used by itself to certify answer
support.

\paragraph{Metric disagreement is central.}
The multi-metric result is not an appendix detail. A 1.1-point DeBERTa effect,
a 3.2-point second-NLI effect, and a 14.0-point local RAGAS-style effect on
the same generations lead to different scientific narratives. Until RAG
faithfulness metrics are better calibrated against human judgments, papers
should report metric disagreement rather than hide it behind a preferred
scalar. In this setting, metric disagreement is part of the estimand, not a
reporting nuisance.

\paragraph{Deployment claims need cost surfaces.}
The baseline audit rejects a single method ranking. RAPTOR-2L narrowly leads
SQuAD faithfulness in the full head-to-head table, CRAG wins HotpotQA on
faithfulness, hallucination rate, and latency, and Self-RAG underperforms in
this matched short-answer harness. Offline indexing cost and p99 latency are
part of the result, not engineering afterthoughts. Deployment claims require
dataset-specific behavior, mean latency, p99 latency, and offline indexing
cost alongside faithfulness.

\paragraph{Implications for RAG evaluation.}
The common RAG reporting pattern, "retrieval metric up, faithfulness metric up,
therefore method works," is insufficient. We recommend that research papers
report at least:
\begin{enumerate}[leftmargin=1.6em,itemsep=2pt,topsep=2pt]
  \item query-passage similarity and retrieval-set coherence;
  \item paired query-level contrasts with bootstrap intervals;
  \item multiple faithfulness metrics;
  \item threshold transfer across datasets;
  \item hallucination rates with Wilson intervals;
  \item mean latency, p99 latency, and offline indexing cost.
\end{enumerate}
